
\subsubsection{Weight-Space Learning}

Kofinas et al. (2024) demonstrated architecture classification using graph neural networks that process computational graphs constructed from weight tensors. Their approach established that weight-space learning is possible but imposed significant complexity barriers through graph construction and 50+ hours of implementation effort. The resulting GNN classifiers operate as black boxes with no interpretability into which weight properties drive decisions. This work demonstrates that for architecture family classification, complex neural approaches are unnecessary---simple statistical features achieve 88.89\% accuracy with $100\times$ faster extraction and full interpretability. Direct accuracy comparison is infeasible due to different datasets and task scopes.

Unterthiner et al. (2020) explored predicting hyperparameters from weight distributions using histogram-based features, demonstrating that statistical summaries of weights contain predictive information. Their focus was training hyperparameters rather than architectural families, and features required analyzing weight value distributions. This work extends that direction by showing that structural metadata alone (layer names, tensor shapes) suffices for architecture classification without examining weight values.

\subsubsection{Normalization Layer Analysis}

Zhang \& Abdulla (2023) analyzed BatchNorm kernel weights for extracting architectural information but required forward passes to compute BatchNorm running statistics, necessitating model instantiation and GPU resources. Chun (2026) provided theoretical foundations demonstrating that LayerNorm and BatchNorm impose fundamentally different geometric constraints---LayerNorm reduces linear layer condition numbers by factors proportional to feature dimensionality, while BatchNorm enforces spatial normalization suited for image data. This work operationalizes these predictions by exploiting that normalization layer choice reflects architectural paradigm and is directly observable through layer name inspection, confirming CNNs exclusively use BatchNorm (0\% violation) and Transformers predominantly use LayerNorm (14.29\% violation) without requiring forward passes.

\subsubsection{Architectural Fingerprinting}

Fang et al. (2024) empirically observed that heterogeneous architectural structures exhibit diverged parameter importance distributions---convolutional layers have many small-magnitude parameters while attention layers have fewer large-magnitude parameters. This work operationalizes that observation by defining parameter-mass ratio R = conv\_params / (conv\_params + linear\_params), demonstrating exceptionally strong inter-family separation (Cohen's d = 3.202, p < 0.001).

Raghu et al. (2021) visualized CNN and Transformer representations through canonical correlation analysis, demonstrating fundamentally different feature hierarchies. Their work required running inference to extract representations. This approach identifies architectural differences purely from checkpoint structure, suggesting behavioral differences are consequences of structural differences observable without execution.

\subsubsection{Positioning}

Prior work established that (1) weight-space learning is possible via complex architectures (Kofinas et al.), (2) checkpoint-based statistical features extract predictive information (Unterthiner et al.), (3) normalization layers impose distinct geometric constraints (Chun), and (4) heterogeneous structures exhibit diverged parameter distributions (Fang). This contribution synthesizes these insights to demonstrate that simple structural metadata features achieve architecture family classification without the complexity of (1), using checkpoint-only structural metadata (2) without weight value inspection, exploiting (3) and (4) to guide feature design. This is the first demonstration of architecture family classification using structural metadata alone without forward passes, weight distributions, or graph construction.
